% BEGIN COMPONENT SGA2-VIII-CUM-C01
\makeatletter
\providecommand{\sgareaderlabel}[2]{%
  \begingroup
  \def\@currentlabel{#2}%
  \phantomsection\label{#1}%
  \endgroup}
\makeatother

\chapter{The Finiteness Theorem}
\label{VIII}

\section*{1. A Biduality Spectral Sequence\footnote{A reader familiar with
Verdier's language of derived categories will recognize the spectral sequence
associated with a \emph{biduality isomorphism}. Cf. SGA~6, I.}}
\sgareaderlabel{VIII.1}{1}

Let us state the result we wish to establish:

\begin{sgastatement}{Proposition 1.1}\sgareaderlabel{VIII.1.1}{1.1}
Let $A$ be a Noetherian ring and let $I$ be an ideal of $A$. Put
$X=\operatorname{Spec}(A)$ and $Y=V(I)$. Let $M$ be a finitely generated
$A$-module of finite projective dimension. Let
$\mathcal F=\widetilde M$ be the $\OX$-module associated with $M$.
\begin{enumerate}
\item[\textup{1)}] There is a spectral sequence
\[
  \Hloc_Y(X,\mathcal F)\Longleftarrow
  \Ext_Y^p(\Ext^{-q}(M,A),A).
\]
\item[\textup{2)}] There is a spectral sequence
\[
  \SheafH_Y(X,\mathcal F)\Longleftarrow
  \SheafExt_Y^p(\SheafExt^{-q}(\mathcal F,\OX),\OX).
\]
\end{enumerate}
\end{sgastatement}
% END COMPONENT SGA2-VIII-CUM-C01

% VERIFIED SOURCE SEAM: French line 2526 is blank and excluded.

% INTEGRATION CONTROL BEFORE SGA2-VIII-CUM-C02
\let\cF\uF

% BEGIN COMPONENT SGA2-VIII-CUM-C02
Of course, 2) follows from 1) by observing that, if $M$ and $N$ are two
finitely generated $A$-modules and if we put $\mathcal F=\widetilde M$ and
$\mathcal G=\widetilde N$, then there are isomorphisms
\begin{align*}
  \SheafH_Y(\mathcal F)
    &\simeq \widetilde{\Hloc_Y(X,\mathcal F)}, \\
  \SheafExt_Y(\mathcal F,\mathcal G)
    &\simeq \widetilde{\Ext_Y(\mathcal F,\mathcal G)}, \\
  \SheafExt_{\OX}(\mathcal F,\mathcal G)
    &\simeq \widetilde{\Ext_A(M,N)}.
\end{align*}

Let $\CatC$ be the category of $A$-modules, and let $\Ab$ be the category of
abelian groups. Let $\cF$ be the functor
\[
  \begin{aligned}
    \cF\colon \CatC &\longrightarrow \Ab, \\
    M &\longmapsto \Gamma_Y(\widetilde M).
  \end{aligned}
\]
We know from Expos\'e~II that there is an isomorphism of $\partial$-functors
\[
  \Hloc_Y^*(X,\widetilde M)\simeq \R^*\cF(M).
\]

Furthermore, let $\Ext_Y^*$ denote the right derived functors, in the second
argument, of
\[
  \cF\circ\operatorname{Hom}\colon \CatC^{\circ}\times\CatC\longrightarrow\Ab.
\]
We know from Expos\'e~VI that there is an isomorphism of $\partial$-functors
\[
  \Ext_Y^*(M,N)\simeq \Ext_Y^*(\widetilde M,\widetilde N).
\]

Finally, let us record the following result from Expos\'e~VI: if $C$ is an
injective $A$-module and $N$ is a finitely generated $A$-module, the sheaf
\[
  \SheafHom(\widetilde N,\widetilde C)
    \simeq \widetilde{\operatorname{Hom}(N,C)}
\]
is flasque; hence
\[
  \R^1\cF(\operatorname{Hom}(N,C))=0.
\]
% END COMPONENT SGA2-VIII-CUM-C02

% VERIFIED SOURCE SEAM: French line 2552 is blank and excluded.

% BEGIN COMPONENT SGA2-VIII-CUM-C03
It remains for us to prove the following result.

\medskip
\sgareaderlabel{VIII.1.2}{1.2}\noindent\textbf{Lemma 1.2.} --- Let $A$ be a Noetherian ring and let $\CatC$
be the category of $A$-modules. Let
$\cF\colon\CatC\longrightarrow\Ab$ be a left exact additive functor such that,
for every finitely generated $A$-module $N$ and every injective $A$-module
$C$, one has
\[
  \R^1\cF(\operatorname{Hom}(N,C))=0.
\]
Let $M$ be a finitely generated $A$-module of finite projective dimension.
There is a spectral sequence
\[
  \R^\ast\cF(M)
    \Longleftarrow
  \Ext_{\cF}^{p}(\Ext^{-q}(M,A),A),
\]
where $\Ext_{\cF}^{p}$ denotes the $p$-th right derived functor of
$\cF\circ\operatorname{Hom}$.
% END COMPONENT SGA2-VIII-CUM-C03

% VERIFIED SOURCE SEAM: French line 2560 is blank and excluded.

% BEGIN COMPONENT SGA2-VIII-CUM-C04
\emph{We shall consider only complexes whose differential has degree $+1$.}
By the hypothesis on $M$, there exists a projective resolution of $M$ of
finite length,
\[
  u\colon L^\bullet\longrightarrow M,
\]
where, moreover, the $L^p$ are finitely generated modules and $L^p=0$ if
$p\notin[-n,0]$. On the other hand, let
$v\colon M\longrightarrow I^\bullet$ be an injective resolution of $M$. I
claim that
\begin{equation}
  v\circ u\colon L^\bullet\longrightarrow I^\bullet
  \tag{1.1}\label{eq:VIII.1.1}
\end{equation}
is an injective resolution of $L^\bullet$. We must specify what this means.

\medskip
\sgareaderlabel{VIII.1.3}{1.3}\noindent\textbf{Definition 1.3.} --- Let $X^\bullet$ be a complex of
$A$-modules. An \emph{injective resolution} of $X^\bullet$ is a homomorphism
of complexes
\[
  x\colon X^\bullet\longrightarrow CX^\bullet
\]
such that $CX^p$ is injective for every $p\in\mathbb Z$, and $x$ induces an
isomorphism in homology.
% END COMPONENT SGA2-VIII-CUM-C04

% VERIFIED SOURCE SEAM: French line 2572 is blank and excluded.

% BEGIN COMPONENT SGA2-VIII-CUM-C05
\sgareaderlabel{VIII.1.4}{1.4}\noindent\textbf{Proposition 1.4.} --- Every left-bounded complex, i.e., every
complex for which there exists $q\in\ZZ$ such that $X^p=0$ for $p<q$, admits
an injective resolution. Moreover, suppose that
$u\colon X^\bullet\longrightarrow Y^\bullet$ is a homomorphism of complexes
(both left-bounded), and that
$x\colon X^\bullet\longrightarrow CX^\bullet$ and
$y\colon Y^\bullet\longrightarrow CY^\bullet$ are injective resolutions of
$X^\bullet$ and $Y^\bullet$. Then there exists a homomorphism of complexes
\[
  Cu\colon CX^\bullet\longrightarrow CY^\bullet,
\]
unique up to homotopy, such that the diagram
\[
\begin{CD}
X^\bullet @>{x}>> CX^\bullet \\
@V{u}VV @VV{Cu}V \\
Y^\bullet @>{y}>> CY^\bullet
\end{CD}
\]
commutes up to homotopy.
% END COMPONENT SGA2-VIII-CUM-C05

% VERIFIED SOURCE SEAM: French line 2582 is blank and excluded.

% BEGIN COMPONENT SGA2-VIII-CUM-C06
The proof is left to the reader.\footnote{See also H.~Cartan and
S.~Eilenberg, \emph{Homological Algebra}, Princeton Mathematical Series,
vol.~19, Princeton University Press, 1956.}

Let us recall a notation introduced in Expos\'e V.

\medskip
\noindent\textbf{Notation.} --- Let $X^\bullet$ and $Y^\bullet$ be two
complexes. We denote by $\operatorname{Hom}^\bullet(X^\bullet,Y^\bullet)$ the
\emph{simple complex} whose component of degree $n$ is
\[
  (\operatorname{Hom}^\bullet(X^\bullet,Y^\bullet))^n
  =\prod_{-p+q=n}\operatorname{Hom}(X^p,Y^q),
\]
also denoted by $\operatorname{Hom}^n(X^\bullet,Y^\bullet)$, and whose
differential is given by
\begin{gather*}
  \partial_n\colon
  \operatorname{Hom}^n(X^\bullet,Y^\bullet)
  \longrightarrow
  \operatorname{Hom}^{n+1}(X^\bullet,Y^\bullet),\\
  \partial_n=d'+(-1)^{n+1}d'',
\end{gather*}
where $d'$ and $d''$ are the differentials (of degree $+1$) induced by those
of $X^\bullet$ and $Y^\bullet$.
% END COMPONENT SGA2-VIII-CUM-C06

% VERIFIED SOURCE SEAM: French line 2596 is blank and excluded.

% INTEGRATION CONTROL BEFORE SGA2-VIII-CUM-C07
\renewcommand{\theequation}{1.\arabic{equation}}\setcounter{equation}{1}

% BEGIN COMPONENT SGA2-VIII-CUM-C07
Let $A^\bullet$ then be the complex defined by $A^p=0$ if $p\ne 0$ and
$A^0=A$. Let
\[
  a\colon A^\bullet\longrightarrow CA^\bullet
\]
be an injective resolution of $A^\bullet$. Consider the double complex
\begin{equation}\label{eq:VIII.1.2}
  Q^{p,q}=\operatorname{Hom}(\operatorname{Hom}(L^{-q},A),CA^p).
\end{equation}
The first spectral sequence of the bicomplex $\cF Q^{\bullet\bullet}$ will
give the conclusion of \hyperref[VIII.1.2]{Lemma~1.2}.

Set
\begin{equation}\label{eq:VIII.1.3}
  {L'}^\bullet=\operatorname{Hom}^\bullet(L^\bullet,A^\bullet),
\end{equation}
and
\begin{equation}\label{eq:VIII.1.4}
  P^\bullet=\operatorname{Hom}^\bullet({L'}^\bullet,CA^\bullet).
\end{equation}
% END COMPONENT SGA2-VIII-CUM-C07

% VERIFIED SOURCE SEAM: French line 2610 is blank and excluded.

% INTEGRATION CONTROL BEFORE SGA2-VIII-CUM-C08
\renewcommand{\theequation}{1.\arabic{equation}}\setcounter{equation}{4}

% BEGIN COMPONENT SGA2-VIII-CUM-C08
One sees easily that $P^\bullet$ is the simple complex associated with
$Q^{\bullet\bullet}$. Let us compute the abutment of the spectral sequence,
that is, the homology of $\uF P^\bullet$. To this end, using the fact that
$L^\bullet$ is projective of finite type in every degree, one proves that
\[
  L^\bullet \cong \operatorname{Hom}^\bullet({L'}^\bullet,A^\bullet).
\]
From the homomorphism $a\colon A^\bullet\to CA^\bullet$, one obtains a
homomorphism
\[
  b\colon \operatorname{Hom}^\bullet({L'}^\bullet,A^\bullet)
  \longrightarrow \operatorname{Hom}^\bullet({L'}^\bullet,CA^\bullet),
\]
or equivalently a homomorphism
\begin{equation}\label{eq:VIII.1.5}
  c\colon L^\bullet\longrightarrow P^\bullet.
\end{equation}

This being said, it is easy to see, using the fact that ${L'}^\bullet$ is
projective of finite type in every degree and bounded on the left, that
\hyperref[eq:VIII.1.5]{(1.5)} is an injective resolution of $L^\bullet$. Using \hyperref[VIII.1.4]{Proposition~1.4}, one
concludes that $P^\bullet$ is homotopy equivalent to $I^\bullet$, where
$I^\bullet$ is the injective resolution of $M$ introduced above in \hyperref[eq:VIII.1.1]{(1.1)}.
It follows that the \emph{abutment} of the first spectral sequence of the
double complex $\uF Q^{\bullet\bullet}$, which is
$H^*(\uF P^\bullet)$, \emph{is isomorphic} to $R^*\uF(M)$.
% END COMPONENT SGA2-VIII-CUM-C08

% VERIFIED SOURCE SEAM: French line 2617 is blank and excluded.

% BEGIN COMPONENT SGA2-VIII-CUM-C09
The initial term of the first spectral sequence of the bicomplex
$\cF Q^{\bullet\bullet}$ is
\[
  E_2^{p,q}
  ={}'\!H^p({}''\!H^q(\cF Q^{\bullet\bullet})).
\]

For every $p\in\ZZ$, $CA^p$ is injective. By the hypothesis on $\cF$, the
functor (restricted to the category of finitely generated modules)
\[
  N\longmapsto \cF\operatorname{Hom}(N,CA^p)
\]
is exact. Hence one obtains isomorphisms
\[
  {}''\!H^q(\cF\operatorname{Hom}({L'}^\bullet,CA^p))
  \cong
  \cF\operatorname{Hom}(H^{-q}({L'}^\bullet),CA^p).
\]
By the definition of $\operatorname{Ext}^{\ast}_{\cF}$ as the right-derived functor
of $\cF\circ\operatorname{Hom}$, one therefore obtains isomorphisms
\[
  E_2^{p,q}
  \cong
  \operatorname{Ext}^{p}_{\cF}(H^{-q}({L'}^\bullet),A).
\]

\noindent\emph{Source note.} The French TeX and printed display have
$H^q({L'}^\bullet)$ in the last formula. The preceding displayed
isomorphism, the statement of \hyperref[VIII.1.2]{Lemma~1.2}, and the identity below require
$H^{-q}({L'}^\bullet)$. This target therefore uses $H^{-q}$ as a transparent
editorial emendation. It is not a restoration from an independent original
scan; the French authority remains unchanged pending editorial action.

Now
\[
  {L'}^\bullet=\operatorname{Hom}^\bullet(L^\bullet,A^\bullet),
\]
where $L^\bullet$ is a projective resolution of $M$; hence there are
isomorphisms
\[
  \operatorname{Ext}^{q}(M,A)\cong H^q({L'}^\bullet).
\]
Substituting $-q$ in this last identity gives the conclusion.\hfill$\square$
% END COMPONENT SGA2-VIII-CUM-C09

% VERIFIED SOURCE SEAM: French line 2639 is blank and excluded.

% INTEGRATION CONTROL BEFORE SGA2-VIII-CUM-C10
\let\cF\calF\renewcommand{\theequation}{\thesection.\arabic{equation}}

% BEGIN COMPONENT SGA2-VIII-CUM-C10
\setcounter{section}{1}
\section{The Finiteness Theorem}\label{VIII.2}

\begin{theorem}\label{VIII.2.1}
\markedfootnote{(1)}{\textit{Editor's note.} For an analogous statement, but
in a somewhat more general situation, see Mme Raynaud (Raynaud M., ``Th\'eor\`emes de
Lefschetz en cohomologie des faisceaux coh\'erents et en cohomologie
\'etale. Application au groupe fondamental,'' \emph{Ann. Sci. \'Ec. Norm.
Sup.} (4) \textbf{7} (1974), pp.~29--52, Proposition~II.2.1).}%
\kern-.18em\textbf{.}\enspace Let $X$ be a locally Noetherian prescheme, let $Y$ be a closed subset of $X$,
and let $\cF$ be a coherent $\OX$-module. Suppose that $X$ is locally
embeddable in a regular prescheme\markedfootnote{*}{This condition may be
generalized to the hypothesis that a ``dualizing complex,'' in the sense
defined in R.~Hartshorne, \emph{Residues and Duality} (cited in the note to
Expos\'e~IV, p.~46), exists locally on $X$.}. Let $i\in\ZZ$. Suppose that:
\begin{enumerate}[label=\textup{\alph*)},leftmargin=2.2em,itemsep=0.5\baselineskip]
\item for every $x\in U=X-Y$, one has
\begin{equation*}
  H^{i-c(x)}(\cF_x)=0,
\end{equation*}
where we have set\markedfootnote{(2)}{\textit{Editor's note.} As in
Expos\'e~V, $H^{\ast}(\cF_x)$ denotes the local cohomology
$H^{\ast}_{\mathfrak m_x}(\cF_x)$.}:
\begin{equation}\label{eq:VIII.2.1}
  c(x)=\operatorname{codim}
  (\overline{\{x\}}\cap Y,\overline{\{x\}}).
\end{equation}
Then:
\item $\mathcal H^i_Y(\cF)$ is coherent.
\end{enumerate}
\end{theorem}
% END COMPONENT SGA2-VIII-CUM-C10

% VERIFIED SOURCE SEAM: French line 2660 is blank and excluded.

% BEGIN COMPONENT SGA2-VIII-CUM-C11
\setcounter{section}{2}
\setcounter{sourcecorollary}{1}
\begin{sourcecorollary}{\textit{Editor's note.} Strictly speaking, this is a
corollary of the proof that follows, not of the statement. The implication
\textup{c)}$\Rightarrow$\textup{a)} is tautological. The converse is not, but
follows from the proof. More precisely, as below, one covers $X$ by open
subsets embeddable in regular schemes; as explained below, this permits one to
reduce to the case where $X=\operatorname{Spec}(A)$ is affine and regular and
$\cF=\widetilde M$, where $M$ is an $A$-module of finite projective dimension.
In this case it is shown that conditions \textup{a)} and \textup{c)} are
equivalent to the dual conditions \textup{a$'$)} and \textup{c$'$)}. One then
shows that \textup{c$'$)} implies condition \textup{d)} (see below), which in
turn implies \textup{a$'$)}. See the discussion following~\hyperref[VIII.2.4]{2.4}.}\label{VIII.2.2}
Under the hypotheses of the preceding theorem, condition \textup{a)} is
equivalent to:
\begin{enumerate}[label=\textup{\alph*)},start=3,leftmargin=2.2em]
\item for every $x\in U$ such that $c(x)=1$, one has
  $H^{i-1}(\cF_x)=0$.
\end{enumerate}
\end{sourcecorollary}
% END COMPONENT SGA2-VIII-CUM-C11

% VERIFIED SOURCE SEAM: French line 2669 is blank and excluded.

% BEGIN COMPONENT SGA2-VIII-CUM-C12
\setcounter{section}{2}
\setcounter{sourcecorollary}{2}
\begin{sourcecorollary}\label{VIII.2.3}
Let $X$ be a locally Noetherian prescheme that is locally embeddable in a
regular prescheme, let $Y$ be a closed subset of $X$, let $\cF$ be a coherent
$\mathcal O_X$-module, and let $n$ be an integer. The following conditions are
equivalent:
\begin{enumerate}[label=\textup{(\roman*)},leftmargin=2.8em]
\item for every $x\in U$, one has
  $\prof \cF_x>n-c(x)$;
\item for every $x\in U$ such that $c(x)=1$, one has
  $\prof \cF_x\geq n$;
\item for every $i\in\mathbb Z$, $\mathcal H^i_Y(\cF)$ is coherent if
  $i\leq n$;
\item $R^i i_*(\cF|_U)$ is coherent for $i<n$%
  \markedfootnote{(4)}{\textit{Editor's note.} This condition appeared only
  in the body of the proof, but not in the statement of the corollary; since
  it is used in \hyperref[VIII.3]{Section~3}, it has been added.}.
\end{enumerate}
\end{sourcecorollary}
% END COMPONENT SGA2-VIII-CUM-C12

% VERIFIED SOURCE SEAM: French line 2682 is blank and excluded.

% BEGIN COMPONENT SGA2-VIII-CUM-C13
\medskip
Assume these results established when $X$ is the spectrum of a regular
Noetherian ring $A$ and when $\cF$ is the sheaf associated with an $A$-module
of finite projective dimension.

First observe that, if $(X_j)_{j\in J}$ is an open covering of $X$ by open
subsets embeddable in a regular scheme, each of the conditions above is
equivalent to the conjunction of the analogous conditions obtained by
replacing $X$ by $X_j$, $Y$ by $Y_j=Y\cap X_j$, and $\cF$ by
$\cF|_{X_j}$. Indeed, only the conditions involving $c(x)$ can present a
difficulty. Let $x\in U$. If $x\in X_j$, set
\[
  c_j(x)=\codim(X_j\cap\overline{\{x\}}\cap Y,
                X_j\cap\overline{\{x\}}).
\]
Necessarily $c_j(x)\geq c(x)$. Let $y\in\overline{\{x\}}\cap Y$ be a point
that ``realizes the codimension,'' that is, such that
$c(x)=\dim\mathcal O_{\overline{\{x\}},y}$, and let $X_j$ be a member of
the covering such that $y\in X_j$. Then $x\in X_j$, and hence
$c_j(x)=c(x)$. It follows that condition~(a) for the $X_j$ implies
condition~(a) for $X$. At this point we have only a partial converse: condition
(a) for $X$ implies condition~(a) for those $X_j$ such that
$c(x)=c_j(x)$, which suffices for our purposes.%
\markedfootnote{(5)}{\textit{Editor's note.} In fact, condition~(a) for $X$
implies condition~(a) for every $X_j$, as asserted in the original text, but
establishing this requires reading the proof below in detail. At this point
the implication does not appear formal. For a property (a), (b), or (c),
append a subscript $J$ to denote its conjunction over the $X_j$. The proof
\textit{below} establishes $(c_J)\Rightarrow(a_J)$ (this is the chain of implications
$(c')\Rightarrow(d)\Rightarrow(a')$). But tautologically
$(a)\Rightarrow(c)$, and $(c)\Longleftrightarrow(c_J)$; hence
$(a)\Rightarrow(a_J)$.}
% END COMPONENT SGA2-VIII-CUM-C13

% VERIFIED SOURCE SEAM: French line 2714 is blank and excluded.

% INTEGRATION CONTROL BEFORE SGA2-VIII-CUM-C14
\setcounter{section}{2}\setcounter{lemma}{3}

% BEGIN COMPONENT SGA2-VIII-CUM-C14
\medskip
Choose a covering of $X$ by open subsets embeddable in a regular prescheme.
Applying the preceding observation, we see that we may assume that $X$ is
closed in a regular $X'$. The reduction to $X'$ is then immediate.

We may therefore assume that $X$ is regular, and even affine, by covering
$X$ with affine open subsets. That we may assume that $\cF=\widetilde M$, where
$M$ has finite projective dimension, will follow from the next lemma.

\begin{lemma}\label{VIII.2.4}
Let $X$ be a regular Noetherian prescheme. Let $\cF$ be a coherent
$\mathcal O_X$-module. The function that assigns to each $x\in X$ the
projective dimension of $\cF_x$ is bounded above.
\end{lemma}
% END COMPONENT SGA2-VIII-CUM-C14

% VERIFIED SOURCE SEAM: French line 2722 is blank and excluded.

% BEGIN COMPONENT SGA2-VIII-CUM-C15
\setcounter{equation}{1}
\medskip
Indeed, let $x\in X$, and let $U$ be an affine open neighborhood of $x$.
Let $L^\bullet$ be a projective resolution of the module $\cF(U)$, where the
$L^i$ are finitely generated. By hypothesis, the ring $\cO_{X,x}$ is regular,
so the projective dimension of $\cF_x$ is finite; denote this integer by $d$.
Set
\[
  K=\ker(L^{-d}\longrightarrow L^{-d+1}).
\]
The module $K_x$ is free, since $d$ is the projective dimension of $\cF_x$
([M], Ch.~VI, Prop.~2.1). By (EGA $0_{\mathrm I}$, 5.4.1 Errata), it follows
that the $\cO_U$-module $\widetilde K$ is free on a neighborhood $U'$ of $x$,
with $U'\subset U$. Choose $f\in\cO_X(U)$ such that
$x\in D(f)\subset U'$. We then have a projective resolution of $M_f$, where
$M=\cF(U)$:
\[
  0\longrightarrow K_f\longrightarrow (L^{-d})_f\longrightarrow
  \cdots\longrightarrow M_f\longrightarrow 0.
\]

\noindent\emph{Source note.} After choosing only $f$, French line 2729 and
the compiled page print $M_{f'}$, whereas the displayed resolution ends in
$M_f$. Moreover, the corrected kernel is a submodule of $L^{-d}$, but French
line 2730 prints $(L^{d-1})_f$. The target uses $M_f$ and $(L^{-d})_f$ as
transparent editorial emendations forced by the chosen localization, the
kernel definition, and the displayed resolution. These are not restorations
from an independent original scan; the French authority remains unchanged
pending editorial review.

This proves that the function under consideration is upper semicontinuous.
Since $X$ is quasi-compact, the conclusion follows.
% END COMPONENT SGA2-VIII-CUM-C15

% VERIFIED SOURCE SEAM: French line 2732 is blank and excluded.

% BEGIN COMPONENT SGA2-VIII-CUM-C16
\medskip
We henceforth assume that $X$ is affine, Noetherian, and regular, and that
$\cF=\widetilde M$, where $M$ is a finitely generated $A$-module, necessarily
of finite projective dimension. We shall proceed in several steps. First, we
find a condition~(d), equivalent to~(a), and prove that it is also equivalent
to~(c). Then, using the spectral sequence of the preceding subsection, we
prove $(d)\Rightarrow(b)$. It remains to prove
$(\mathrm{iii})\Rightarrow(\mathrm{ii})$; indeed,
$(\mathrm{i})\Leftrightarrow(\mathrm{ii})\Rightarrow(\mathrm{iii})$ follows
immediately from $(a)\Leftrightarrow(c)\Rightarrow(b)$.

Let $x\in U$. By hypothesis, $\cO_{X,x}$ is a regular local ring. If $D$
denotes the dualizing functor for $\cO_{X,x}$, it follows from (\hyperref[V.2.1]{V, 2.1)} that
\[
  D\,H^{i-c(x)}(\cF_x)\simeq
  \operatorname{Ext}^{d(x)-i}_{\cO_{X,x}}(\cF_x,\cO_{X,x}),
\]
where
\begin{equation}\label{eq:VIII.2.2}
  d(x)=\dim\cO_{X,x}+c(x)
      =\dim\cO_{X,x}
       +\codim(\overline{\{x\}}\cap Y,\overline{\{x\}}).
  \tag{2.2}
\end{equation}
Since $X$ is Noetherian and $\cF$ is coherent,
\begin{equation}\label{eq:VIII.2.3}
  D\,H^{i-c(x)}(\cF_x)\simeq
  (\sExt^{d(x)-i}_{\cO_X}(\cF,\cO_X))_x.
  \tag{2.3}
\end{equation}
Moreover, a module vanishes if and only if its dual vanishes (cf. the
editor's note~(4) on p.~54). For every $q\in\mathbb Z$, set
\begin{equation}\label{eq:VIII.2.4}
  \begin{cases}
    S_q=\Supp\sExt^q_{\cO_X}(\cF,\cO_X),\\
    S'_q=S_q\cap U,\quad (U=X-Y),\\
    Z_q=\overline{S'_q}\cap Y.
  \end{cases}
  \tag{2.4}
\end{equation}
% END COMPONENT SGA2-VIII-CUM-C16

% BEGIN COMPONENT SGA2-VIII-CUM-C17
\medskip
It follows from \hyperref[eq:VIII.2.3]{formula~(2.3)} that conditions~(a) and~(c) are respectively
equivalent to:

\medskip
\noindent\textup{(a$'$)} for every $q\in\mathbb Z$ and every $x\in S'_q$,
one has $q+i\ne d(x)$;

\smallskip
\noindent\textup{(c$'$)} for every $q\in\mathbb Z$ and every $x\in S'_q$,
if $c(x)=1$, then $q+i\ne d(x)$.

\medskip
\noindent Here is condition~(d), promised above:

\medskip
\noindent\textup{(d)} for every $q\in\mathbb Z$ and every $y\in Z_q$,
one has $q+i\ne\dim\cO_{X,y}$.

\medskip
\noindent These conditions are equivalent.

\smallskip
\noindent\textup{(a$'$)} $\Rightarrow$ \textup{(c$'$)} is recorded for
completeness.

\smallskip
\noindent\textup{(d)} $\Rightarrow$ \textup{(a$'$)}. Indeed, let
$q\in\mathbb Z$ and $x\in S'_q$. Choose a point
$y\in\overline{\{x\}}\cap Y$ that realizes the
codimension,\markedfootnote{(6)}{\textit{Editor's note.} Throughout what
follows, closures of points are endowed with the reduced structure.} that is,
such that
\begin{equation}\label{eq:VIII.2.5}
  \dim\cO_{\overline{\{x\}},y}
  =\codim(\overline{\{x\}}\cap Y,\overline{\{x\}})=c(x).
  \tag{2.5}
\end{equation}
Since $X$ is regular at $y$, it follows that
\begin{equation}\label{eq:VIII.2.6}
  \dim\cO_{X,y}=d(x)\qquad\text{(cf.~\hyperref[eq:VIII.2.2]{(2.2)}).}
  \tag{2.6}
\end{equation}
But $y\in\overline{\{x\}}$, and hence $y\in Z_q$; the conclusion follows.

\smallskip
\noindent\textup{(c$'$)} $\Rightarrow$ \textup{(d)}. Let $q\in\mathbb Z$
and $y\in Z_q$. Assume provisionally that there exists $x\in S'_q$ such
that
\[
  y\in\overline{\{x\}}\qquad\text{and}\qquad
  \dim\cO_{\overline{\{x\}},y}=1.
\]
(One also says that $y$ is an immediate specialization of $x$.) It follows
that $c(x)=1$, since $y$ realizes the codimension of
$\overline{\{x\}}\cap Y$ in $\overline{\{x\}}$, while $x\notin Y$.
Condition~\textup{(c$'$)} therefore gives
\[
  q+i\ne d(x).
\]
Whence the conclusion, upon observing that $d(x)=\dim\cO_{X,y}$ by~\hyperref[eq:VIII.2.6]{(2.6)}.
The assertion provisionally assumed above is the content of the following
lemma.

\medskip
\sgareaderlabel{VIII.2.5}{2.5}\noindent\textbf{Lemma 2.5.} Let $X$ be a locally Noetherian prescheme and
$Y$ a closed subset of $X$. Set $U=X-Y$, and suppose that $U$ is dense in
$X$. For every $y\in Y$, there exists $x\in U$ such that $y$ is an immediate
specialization of $x$, that is,
\[
  y\in\overline{\{x\}}\qquad\text{and}\qquad
  \dim\cO_{\overline{\{x\}},y}=1.
\]

We have applied the lemma with $X$ equal to the prescheme
$\overline{S'_q}$ and $Y$ equal to the subset
$Y\cap\overline{S'_q}$.

\medskip
\noindent\textit{Proof of \hyperref[VIII.2.5]{Lemma 2.5}.} There exists $x\in U$ such that
$y\in\overline{\{x\}}$. Choose such an $x$ for which
$r=\dim\cO_{\overline{\{x\}},y}$ is minimal. We must prove that $r=1$.
By the choice of $x$, every
$z\in\Spec(\cO_{\overline{\{x\}},y})$ with $z\ne x$ lies in $Y$.
Consequently, $\{x\}$ is open in
$\Spec(\cO_{\overline{\{x\}},y})$, and the conclusion follows.\hfill$\square$
% END COMPONENT SGA2-VIII-CUM-C17

% VERIFIED SOURCE SEAM: French line 2795 is blank and excluded.

% BEGIN COMPONENT SGA2-VIII-CUM-C18
\medskip
\noindent\textit{The second step is to deduce \textup{(b)} from
\textup{(d)}.}

\medskip
Set
\[
  D(Z_q)=\{\dim\cO_{X,y}\mid y\in Z_q\}.
\]
By~\textup{(d)}, for every $q\in\mathbb Z$ one has
$q+i\notin D(Z_q)$. Applying Proposition~\hyperref[VII.2.3]{VII.2.3}, we see that
\[
  \sExt_Y^{q+i}(\sExt^q(\cF,\cO_X),\cO_X)
  \quad\text{is coherent.}
\]
The initial term of the spectral sequence of the preceding subsection is
\[
  E_2^{p,q}
  =\sExt_Y^p(\sExt^{-q}(\cF,\cO_X),\cO_X).
\]
It follows that $E_2^{p,q}$ is coherent for every $p,q\in\mathbb Z$ such
that $p+q=i$. Now there are only finitely many pairs $(p,q)$ such that
$p+q=i$,\markedfootnote{[S]}{\textit{Source note.} Taken literally for
arbitrary integers, this sentence is false. In context it must be read as
concerning the terms that can contribute, in view of the finite-projective-dimension
setup above. The French authority is unchanged.} and this spectral
sequence converges to $\sH_Y^\ast(\cF)$; whence the conclusion.
% END COMPONENT SGA2-VIII-CUM-C18

% VERIFIED SOURCE SEAM: French line 2806 is blank and excluded.

% BEGIN COMPONENT SGA2-VIII-CUM-C19
\medskip
\noindent\textit{It remains to prove that \textup{(iii)} implies
\textup{(ii)}.} Let
\[
  i\colon U\longrightarrow X
\]
denote the canonical immersion of $U$ into $X$. In view of the exact homology
sequence associated with the closed subset $Y$ (\hyperref[I.2.11]{I, 2.11)}, one sees that
\textup{(iii)} is \textit{equivalent to}:

\medskip
\noindent\textup{(iv)} $R^i i_*(\cF|_U)$ is coherent for $i<n$.
\refstepcounter{sourcecondition}\label{condition}

\medskip
\noindent{\scriptsize\color{gray}French printed p.~94}\quad Indeed, there is
an exact sequence
\[
  0\longrightarrow \sH_Y^0(\cF)\longrightarrow \cF
  \longrightarrow i_*(\cF|_U)\longrightarrow \sH_Y^1(\cF)
  \longrightarrow 0.
\]
Now $\sH_Y^0(\cF)$ is a quasi-coherent subsheaf of the coherent sheaf $\cF$,
and is therefore coherent. Hence $\sH_Y^1(\cF)$ is coherent if and only if
$i_*(\cF|_U)$ is coherent. Moreover, for $p>0$, the exact cohomology sequence
associated with the closed subset $Y$ reduces to isomorphisms
\[
  R^p i_*(\cF|_U)\overset{\approx}{\longrightarrow}\sH_Y^{p+1}(\cF).
\]
% END COMPONENT SGA2-VIII-CUM-C19

% BEGIN COMPONENT SGA2-VIII-CUM-C20
\medskip
\noindent\textit{We shall prove that \textup{(iv)} implies
\textup{(ii)}.} To this end, recall \textup{(ii)}:

\medskip
\noindent\textup{(ii)} For every $x\in U$ such that $c(x)=1$, one has
\[
  \operatorname{prof}\cF_x\geq n.
\]

\noindent We argue by induction on $n$.

\medskip
\noindent If $n=0$, both conditions are vacuous.
% END COMPONENT SGA2-VIII-CUM-C20

% VERIFIED SOURCE SEAM: French line 2827 is blank and excluded.

% BEGIN COMPONENT SGA2-VIII-CUM-C21
\medskip
\noindent If $n=1$, assume that $i_*(\cF|_U)$ is coherent. We argue by
contradiction and suppose that there exists $x\in U$ such that $c(x)=1$ and
\[
  \prof\cF_x=0,
\]
i.e. $x\in\Ass\cF_x$. Let $y\in\overline{\{x\}}\cap Y$ be such that
\[
  \dim\cO_{\overline{\{x\}},y}=1.
\]
Set
\[
  A=\cO_{X,y}\qquad\text{and}\qquad X'=\Spec(A).
\]

\setcounter{equation}{6}
\noindent Perform the base change $v\colon X'\longrightarrow X$, which is
flat:
\begin{equation}
\begin{CD}
U'=X'\times_X U @>{v'}>> U\\
@V{i'}VV @VV{i}V\\
X' @>{v}>> X .
\end{CD}
\tag{2.7}\label{eq:VIII.2.7}
\end{equation}
The morphism $i$ is separated, since it is an immersion, and is of finite
type, since it is an open immersion and $X$ is locally Noetherian. Since the
base change is flat, EGA~III, 1.4.15 gives an isomorphism
\begin{equation}
  v^*(i_*(\cF|_U))
  \approx
  i'_*(v'^*(\cF|_U)).
\tag{2.8}\label{eq:VIII.2.8}
\end{equation}

\noindent{\scriptsize\color{gray}French printed p.~95}\quad Denote by
$\underline{x}$ (respectively, $\underline{y}$) the ideal of $A$
corresponding to $x$ (respectively, $y$).

\medskip
\noindent Set $\cG=v'^*(\cF|_U)$. Then $\cG$ is coherent and
$\underline{x}\in\Ass\cG$, so there exists a monomorphism
\[
  \cO_{\overline{\{x\}}}\longrightarrow\cG;
\]
consequently, $i'_*(\cO_{\overline{\{x\}}}|_{U'})$ is coherent. By the choice
of $y$, $\dim(A/\underline{x})=1$, and therefore the support of
$\cO_{\overline{\{x\}}}$ is precisely
\[
  \overline{\{\underline{x}\}}
  =\{\underline{x}\}\cup\{\underline{y}\},
\]
since $\overline{\{\underline{x}\}}=\Spec(A/\underline{x})$ as a scheme. It
follows that
\[
  (\cO_{\overline{\{x\}}}|_{U'})(U')
  =\Frac(A/\underline{x}),
\]
the field of fractions of $A/\underline{x}$, and
\[
  i'_*(\cO_{\overline{\{x\}}}|_{U'})(X')
  =\Frac(A/\underline{x}).
\]
But $\Frac(A/\underline{x})$ is not a finitely generated $A$-module, because
$\underline{x}$ is different from the maximal ideal of $A$. This is a
contradiction.
% END COMPONENT SGA2-VIII-CUM-C21

% VERIFIED SOURCE SEAM: French line 2855 is blank and excluded.

% BEGIN COMPONENT SGA2-VIII-CUM-C22
\medskip
\noindent Suppose that $n>1$, and assume the result established for every
$n'<n$. By the induction hypothesis, for every $x\in U$ such that $c(x)=1$,
one has $x\notin\Ass\cF_x$. Let $x$ be such a point, and let
$y\in\overline{\{x\}}\cap Y$ be an immediate specialization of $x$, i.e.
\[
  \dim\cO_{\overline{\{x\}},y}=1.
\]
Perform the base change
\[
  v\colon\Spec(\cO_{X,y})\longrightarrow X,
\]
retaining the notation of diagram~\textup{\hyperref[eq:VIII.2.7]{(2.7)}}. Applying EGA~III, 1.4.15,
one obtains isomorphisms
\[
  v^*(R^p i_*(\cF|_U))
  \simeq
  R^p i'_*(v'^*(\cF|_U)),
  \qquad p\in\mathbb Z.
\]

\noindent\emph{Thus we are reduced to the case in which $X$ is the spectrum
of a local ring $A$ in which $\underline{x}$ is a prime ideal of dimension
$1$}, i.e. $\dim(A/x)=1$. Set
\[
  \cF'=\SheafGamma_Y(\cF)
  \qquad\text{and}\qquad
  \cF''=\cF/\cF'.
\]
We have $\cF_x\simeq\cF''_x$ and $y\notin\Ass\cF''$. Moreover,
$\cF'|_U=0$; hence the long exact sequence of the $R^p i_*$ gives
isomorphisms
\[
  R^p i_*(\cF|_U)
  \simeq
  R^p i_*(\cF''|_U),
  \qquad p\in\mathbb Z.
\]

\noindent{\scriptsize\color{gray}French printed p.~96}\quad Since $n>1$,
neither $x$ nor $y$ belongs to $\Ass\cF''$. Now $\underline{x}$ and
$\underline{y}$ are the only prime ideals of $A$ containing $\underline{x}$.
It follows from Expos\'e~\hyperref[III.2.1]{III, 2.1} that there exists an element
$g\in\underline{x}$ that is $M$-regular, where we have set
\[
  \cF=\widetilde M,
  \qquad M=\cF(X).
\]
Thus there is an exact sequence
\[
  0\longrightarrow M\xrightarrow{g'}M\longrightarrow N\longrightarrow0,
\]
in which $g'$ denotes multiplication by $g$ on $M$. From the resulting long
exact sequence one sees that
\[
  R^p i_*(\widetilde N|_U)
  \quad\text{is coherent for }p<n-1.
\]
Therefore, by the induction hypothesis,
\[
  \prof(\widetilde N)_x\geq n-1,
  \qquad\text{and hence}\qquad
  \prof\cF_x\geq n.
\]
\hfill$\square$
% END COMPONENT SGA2-VIII-CUM-C22

% VERIFIED SOURCE SEAM: French line 2873 is blank and excluded.

% BEGIN COMPONENT SGA2-VIII-CUM-C23
\section*{3. Applications}
\sgareaderlabel{VIII.3}{3}

\noindent From these results one deduces a coherence condition for the higher
direct images of a coherent sheaf under a \emph{morphism that is not proper}.

\medskip
\sgareaderlabel{VIII.3.1}{3.1}\noindent\textbf{Theorem 3.1.}---Let $f\colon X\to Y$ be a morphism of
preschemes. Suppose that $Y$ is \emph{locally Noetherian} and that $f$ is
\emph{proper}. Suppose that $X$ is \emph{locally embeddable in a regular
prescheme}. Let $n\in\mathbb Z$. Let $U$ be an open subset of $X$, and let
$\cF$ be a \emph{coherent} $\cO_U$-module. Suppose that, for every $x\in U$
such that
\[
  \codimn(\overline{\{x\}}\cap(X\setminus U),
  \overline{\{x\}})=1,
\]
one has $\prof\cF_x\geq n$. \emph{Then} the $\cO_Y$-modules
\[
  R^p(f\circ g)_*(\cF)
\]
are coherent for $p<n$, where $g$ is the canonical immersion of $U$ into
$X$.

\medskip
\noindent Indeed, there is a Leray spectral sequence whose abutment is
$R^*(f\circ g)_*(\cF)$ and whose initial term is given by
\[
  E_2^{p,q}=R^p f_*(R^q g_*(\cF)).
\]
Moreover, there exists a \emph{coherent} $\cO_X$-module $\cG$ such that
$\cG|_U\simeq\cF$ (EGA~I, 9.4.3). It then follows from the preceding section
that condition~\textup{(iv)} on p.~74 is satisfied
{\scriptsize\color{gray}[French printed p.~97]}, i.e. that
$R^q g_*(\cG|_U)$ is coherent for $q<n$. Applying the finiteness theorem of
EGA~III, 3.2.1 to $f$ and the sheaves $R^q g_*(\cF)$, one finds that
$E_2^{p,q}$ is coherent for $q<n$, whence the conclusion.
% END COMPONENT SGA2-VIII-CUM-C23

% VERIFIED SOURCE SEAM: French line 2887 is blank and excluded.

% BEGIN COMPONENT SGA2-VIII-CUM-C24
\medskip
\sgareaderlabel{VIII.3.2}{3.2}\noindent\textbf{Proposition 3.2.}---Let $X$ be a locally Noetherian
prescheme, locally embeddable in a regular prescheme. Let $U$ be an open
subset of $X$, and let $i\colon U\to X$ be the canonical immersion. Let
$n\in\mathbb Z$. Finally, let $\cF$ be a coherent Cohen--Macaulay
$\cO_U$-module (on $U$!). The following conditions are equivalent:
\begin{enumerate}
\item[(a)] $R^p i_*(\cF)$ is coherent for $p<n$;
\item[(b)] for every irreducible component $S'$ of the closure $\overline S$
of the support $S$ of $\cF$, one has
\[
  \codimn(S'\cap(X\setminus U),S')>n.
\]
\end{enumerate}
% END COMPONENT SGA2-VIII-CUM-C24

% VERIFIED SOURCE SEAM: French line 2900 is blank and excluded.

% BEGIN COMPONENT SGA2-VIII-CUM-C25
\medskip
\noindent\textit{Proof, opening.}---Let $\cG$ be a coherent $\cO_X$-module
such that $\cG|_U\simeq\cF$ (EGA~I, 9.4.3). Applying \hyperref[VIII.2.3]{Corollary~2.3} to
$\cG$, we find that condition \textup{(a)} is equivalent to \textup{(c)}:
\begin{enumerate}
\item[(c)] for every $x\in S$, one has
\[
  \prof\cF_x>n-c(x),
  \qquad
  c(x)=\codimn(\overline{\{x\}}\cap Y,\overline{\{x\}}).
\]
\end{enumerate}
% END COMPONENT SGA2-VIII-CUM-C25

% VERIFIED SOURCE SEAM: French line 2908 is blank and excluded.

% BEGIN COMPONENT SGA2-VIII-CUM-C26
\medskip
\noindent\textit{Proof, \textup{(a)}$\Rightarrow$\textup{(b)}.}---Indeed,
let $S'$ be an irreducible component of $\overline S$, and let $s$ be its
generic point. Since $\cF$ is Cohen--Macaulay, one has
$\prof \cF_s=\dim\cO_{S,s}=0$. Moreover, $\overline{\{s\}}=S'$, and the
conclusion follows.
% END COMPONENT SGA2-VIII-CUM-C26

% VERIFIED SOURCE SEAM: French line 2910 is blank and excluded.

% BEGIN COMPONENT SGA2-VIII-CUM-C27
\medskip
\noindent\textit{Proof, \textup{(b)}$\Rightarrow$\textup{(a)}, opening.}---Let
$x\in S$, and let $S'$ be an irreducible component of $\overline S$ such that
$x\in S'$. Let $s$ be the generic point of $S'$. Since $\cF$ is
Cohen--Macaulay, we know that
\[
  \prof \cF_x=\dim\cO_{\overline{\{s\}},x}.
\]
% END COMPONENT SGA2-VIII-CUM-C27

% BEGIN COMPONENT SGA2-VIII-CUM-C28
\medskip
\noindent\textit{Proof, \textup{(b)}$\Rightarrow$\textup{(a)}, case
split.}---If $c(x)=+\infty$, there is nothing to prove. Otherwise, there exists
$y\in Y\cap\overline{\{x\}}$ such that
\[
  c(x)=\dim\cO_{\overline{\{x\}},y}.
\]
% END COMPONENT SGA2-VIII-CUM-C28

% BEGIN COMPONENT SGA2-VIII-CUM-C29
\medskip
\noindent\textit{Proof, \textup{(b)}$\Rightarrow$\textup{(a)},
conclusion.}---Now $\cO_{X,y}$ is a quotient of a regular local ring by
hypothesis; hence
\[
  \dim\cO_{\overline{\{s\}},y}
  =\dim\cO_{\overline{\{s\}},x}
   +\dim\cO_{\overline{\{x\}},y}>n.
\]
\hfill$\square$
% END COMPONENT SGA2-VIII-CUM-C29

% VERIFIED SOURCE SEAM: French line 2923 is blank and excluded.

% INTEGRATION CONTROL BEFORE SGA2-VIII-CUM-C30
\begingroup

% BEGIN COMPONENT SGA2-VIII-CUM-C30
\renewcommand{\refname}{Bibliography}
\begin{thebibliography}{M}
\bibitem[M]{M}
\textsc{H.~Cartan \& S.~Eilenberg} -- \emph{Homological Algebra},
Princeton Math. Series vol.~19, Princeton University Press, 1956.
\end{thebibliography}
\endgroup
% END COMPONENT SGA2-VIII-CUM-C30
